% supplemental.tex
% Supplemental Material for the PRL companion to claudeFinal.tex.
% Style rules (per AGENT.md, AGENTS.md):
%   - physics prose only; all status / methodology / open notes are % comments
%   - the de Vries number 0.2231 is an output, never an input
%   - integer-fit selection of (p,q,r,n) is explicitly NOT a derivation route here
\documentclass[aps,prd,onecolumn,superscriptaddress,nofootinbib,reprint]{revtex4-2}

\usepackage{amsmath,amssymb,amsthm,mathtools}
\usepackage{bm}
\usepackage{booktabs}
\usepackage{longtable}
\usepackage{graphicx}
\usepackage{hyperref}

\newtheorem{theorem}{Theorem}
\newtheorem{proposition}{Proposition}
\newtheorem{lemma}{Lemma}

\newcommand{\Xp}{X_{+}}
\newcommand{\Xm}{X_{-}}
\newcommand{\half}{\tfrac12}
\newcommand{\sinsq}{\sin^{2}\!\theta_{W}}
\newcommand{\cA}{\mathcal{A}}
\newcommand{\cH}{\mathcal{H}}
\newcommand{\cO}{\mathcal{O}}
\newcommand{\SU}{\mathrm{SU}}
\newcommand{\U}{\mathrm{U}}
\newcommand{\SO}{\mathrm{SO}}
\newcommand{\CP}{\mathbb{CP}}
\newcommand{\AdS}{\mathrm{AdS}}
\newcommand{\dd}{\mathrm{d}}
\newcommand{\R}{\mathbb{R}}
\newcommand{\C}{\mathbb{C}}
\newcommand{\Z}{\mathbb{Z}}

\begin{document}

\title{Supplemental Material:\\
Geometric Origin of the Electroweak Mixing Angle}

\author{[author]}
\affiliation{[affiliation]}

\date{\today}

\maketitle

\tableofcontents

% ============================================================
\section{Normal-form theorem: detailed proof}
\label{app:nf}
% ============================================================

We prove Theorem~1 of the main text.  Let $V_{j}$ be the spin-$j$
representation of $\SU(2)$, with Casimir
$\sum_{a}T_{a}^{(j)}T_{a}^{(j)}=J\,\mathbf{1}$, and let $\cH_{j}$ be the
$2\times 2$ block defined as Eq.~(3) of the Letter, acting on
$V_{j}\oplus\mathrm{Im}\,A_{j}$.

\begin{lemma}[Casimir factorization]
\label{lem:cas}
Let $\{e^{a}\}$ be an orthonormal basis of the cotangent space at any
point of a parallelizable three-manifold with $\SU(2)$ structure
constants.  Then the operator
$A_{j}=\sum_{a}T_{a}^{(j)}\otimes e^{a}$ obeys
$A_{j}^{\dagger}A_{j}=J\,\mathbf{1}_{V_{j}}$ in the orthonormal frame.
\end{lemma}

\begin{proof}
Using $\langle e^{a},e^{b}\rangle=\delta^{ab}$,
\begin{equation*}
  A_{j}^{\dagger}A_{j}
  =\sum_{a,b}T_{a}^{(j)\dagger}T_{b}^{(j)}\,\delta^{ab}
  =\sum_{a}T_{a}^{(j)}T_{a}^{(j)}
  =J\,\mathbf{1}_{V_{j}}.\qedhere
\end{equation*}
\end{proof}

\begin{proof}[Proof of Theorem~1]
Axiom~(i) sets the upper-left entry to zero.  Axiom~(ii) and
Lemma~\ref{lem:cas} fix the modulus of the off-diagonal as
$|\mathcal{O}_{\rm mix}|^{2}=\mu^{2}J$ on the image.  Choosing a phase,
$\mathcal{O}_{\rm mix}=\mu\,A_{j}^{\dagger}$.  Axiom~(iii) then gives
$M_{\rm s}^{2}=-A_{j}A_{j}^{\dagger}$.  Projecting to the image
two-channel subspace (parameterized by $\phi\in V_{j}$ and the image
component of $a$), the block becomes
\begin{equation}
  \mu^{2}
  \begin{pmatrix} 0 & \sqrt{J}\\ \sqrt{J} & -J\end{pmatrix}
  =\mu^{2}\cA(J),
\end{equation}
which is unique by axiom~(iv).  The characteristic equation $X^{2}+JX-J=0$ has
roots
$\Xp(J)=\tfrac12(-J+\sqrt{J^{2}+4J})>0$ and
$\Xm(J)=\tfrac12(-J-\sqrt{J^{2}+4J})<0$, with
$\Xp\Xm=-J$ and $\Xp+\Xm=-J$.
\end{proof}

The structural content of the four axioms is the following.  Axiom~(i)
is gauge masslessness before symmetry breaking.  Axiom~(ii) is the
Casimir factorization, supplied by any first-order
$\SU(2)$-covariant operator on a parallelizable internal three-cycle.
Axiom~(iii) is the single-coupling Schur completion of
Proposition~\ref{prop:c1-app} below.  Axiom~(iv) is the statement that
no further dimensionless coefficient enters.  These four together
force $\cA(J)$ uniquely; the weak mixing angle is then a consequence
of representation theory.

% ============================================================
\section{Gauge--Higgs reduction on \texorpdfstring{$\CP^{1}$}{CP1} and the bare Casimir}
\label{app:gh}
% ============================================================

The off-diagonal $\sqrt{J}$ of $\cA(J)$ is the matrix element of a
specific physical operator built from the gauge--Higgs reduction of the
six-dimensional $\SU(2)$ gauge theory on $M_{4}\times\CP^{1}$.  We
record the reduction in enough detail to fix the normalization.

\subsection{Coset structure and line-bundle modes}

The weak $\CP^{1}\cong\SU(2)/\U(1)_{T_{3}}$ carries a homogeneous round
metric of radius $R$.  In the patch $(\theta,\varphi)$ the
left-invariant generators of $\SU(2)$ act as Killing vectors, and the
$\U(1)_{T_{3}}$ stabilizer of the north pole is the diagonal
$\SU(2)$ subgroup.  The minimal-magnetic-charge background on $\CP^{1}$
is the Hopf line bundle $\cO(1)$, with first Chern number $1$.

The harmonic decomposition on $\CP^{1}$ in the presence of the Hopf
background is the Borel--Weil theorem:
\begin{equation}
  H^{0}\bigl(\CP^{1},\cO(2j)\bigr)\;\cong\;V_{j},\qquad j=0,\half,1,\dots
\label{eq:bw-app}
\end{equation}
i.e.\ the holomorphic sections of $\cO(2j)$ form the spin-$j$
representation of $\SU(2)$ of dimension $2j+1$.  For $j=\half$ this is
the Higgs doublet; for $j=1$ it is the triplet of the same operator;
$j=0$ is the singlet.

\subsection{The six-dimensional gauge action and its 4D content}

Write the six-dimensional $\SU(2)$ gauge field as
$\hat A_{M}=(A_{\mu}(x,y),A_{\alpha}(x,y))$ with $\mu=0,\dots,3$ on
$M_{4}$ and $\alpha=1,2$ on $\CP^{1}$.  Expand the off-diagonal
components $A_{\alpha}^{\pm}=(A_{\alpha}^{1}\pm iA_{\alpha}^{2})/\sqrt 2$
in $\cO(\pm 1)$ harmonics on $\CP^{1}$.  The zero modes of $A_{\alpha}^{+}$
form a doublet $H\in V_{1/2}$ under the surviving four-dimensional
$\SU(2)_{W}$; the zero modes of $A_{\alpha}^{-}$ form the conjugate
doublet $H^{*}$.  Together, $H=(H_{+},H_{0})$ is the Standard-Model
Higgs.

The six-dimensional action is
\begin{equation}
  S_{6}=-\frac{1}{4g_{6}^{2}}\int d^{6}X\sqrt{-g}\,
  F_{MN}^{a}F^{MN,a}.
\label{eq:S6}
\end{equation}
Decomposing $F_{MN}=(F_{\mu\nu},F_{\mu\alpha},F_{\alpha\beta})$ and
integrating $A_{\alpha}$ over the zero modes \eqref{eq:bw-app}, the
4D action acquires (i) the kinetic term of the doublet $H$ from the
$F_{\mu\alpha}^{2}$ piece, (ii) the cross-coupling between $W_{\mu}^{a}$
and $H$ from the cross-strength
$F_{\mu\alpha}=\partial_{\mu}A_{\alpha}-\partial_{\alpha}A_{\mu}+g[A_{\mu},A_{\alpha}]$,
and (iii) a one-loop Hosotani potential for $H$ from the
$F_{\alpha\beta}^{2}$ piece evaluated on the KK tower.

The piece relevant for the off-diagonal of $\cA(J)$ [Eq.~(2) of the Letter] is the
covariant derivative of $H$:
\begin{equation}
  D_{\mu}H\;=\;\partial_{\mu}H\;-\;i g W_{\mu}^{a}\,T_{a}^{(j)}H,
  \qquad j=\half.
\label{eq:DmuH}
\end{equation}
The generators $T_{a}^{(j)}$ in \eqref{eq:DmuH} are the \emph{algebraic}
$\SU(2)$ generators acting on the line-bundle zero modes $V_{j}$, not
the geometric Lie derivative on the cotangent bundle of the internal
three-manifold; this is the bare-Casimir content of gauge--Higgs
unification.

\subsection{The vacuum and the gauge-boson masses}

The vacuum is $\langle H\rangle=v\eta_{0}$ with $\eta_{0}\in V_{1/2}$
the neutral component selected by hypercharge.  Acting with the
algebraic generators,
\begin{equation}
  T_{3}^{(1/2)}\eta_{0}=-\half\,\eta_{0},\qquad
  Y\eta_{0}=+\half\,\eta_{0},\qquad
  Q\eta_{0}=(T_{3}+Y)\eta_{0}=0,
\label{eq:Qeta}
\end{equation}
so the photon direction is exactly massless.  The gauge-boson
mass-squared matrix at the vacuum is
\begin{equation}
  M_{ab}^{2}=g^{2}v^{2}\,\langle\eta_{0}|T_{a}^{(1/2)}T_{b}^{(1/2)}|\eta_{0}\rangle,
\label{eq:Mab}
\end{equation}
whose trace is
\begin{equation}
  \sum_{a}M_{aa}^{2}\;=\;g^{2}v^{2}\,\langle\eta_{0}|\sum_{a}T_{a}T_{a}|\eta_{0}\rangle
  \;=\;g^{2}v^{2}\,J_{1/2}\;=\;\tfrac34\,g^{2}v^{2},
\label{eq:Mab-trace}
\end{equation}
the bare Casimir of $V_{1/2}$.  Diagonalizing \eqref{eq:Mab} with
hypercharge mixing gives the standard $M_{W}^{2}=g^{2}v^{2}/4$ and
$M_{Z}^{2}=(g^{2}+g'^{2})v^{2}/4$.

\subsection{Comparison with the standard textbook reading}
\label{app:standard}

In the textbook gauge--Higgs reading, the Higgs is a Wilson-line
modulus, the gauge bosons get mass through the Higgs mechanism
\eqref{eq:Mab}, and the Higgs mass comes from the one-loop Hosotani
effective potential.  This gives three scales: $M_{W}=gv/2$,
$M_{Z}=\sqrt{g^{2}+g'^{2}}\,v/2$, and $m_{H}=\sqrt{2\lambda}v$ with
$\lambda$ extracted from the loop calculation on the coset.  The
weak mixing angle is the embedding-coupling ratio
$\sin^{2}\theta_{W}=g'^{2}/(g^{2}+g'^{2})$, which on the $SU(3)/U(2)$
coset gives $1/4$~\cite{Manton1979,Fairlie1979}.  Our setting differs:
the weak coset is $\SU(2)/\U(1)\cong\CP^{1}$ and hypercharge comes
from the $\CP^{2}$ factor with its spin\(^{c}\) structure, so the
Manton--Fairlie $1/4$ does not apply.  What does apply is the
mass-squared matrix \eqref{eq:Mab}.

The de Vries reading uses the $(W,\delta H)$ block on the spin-$j$
sector as a single indefinite matrix $\cA(J)$ at one scale $\mu$.
Identifying $M_{W}^{2}=\mu^{2}\Xp(3/4)$ recovers the textbook
$M_{W}^{2}=g^{2}v^{2}/4$ through $\mu^{2}=g^{2}v^{2}/(4\Xp(3/4))$.
The Higgs mass is then read from the negative eigenvalue:
$m_{H}^{2}=\mu^{2}|\Xm(3/4)|$, with no free parameter $\lambda$.  In
ratio form,
\begin{equation}
  \frac{m_{H}^{2}}{M_{W}^{2}}\;=\;\frac{|\Xm(3/4)|}{\Xp(3/4)}\;=\;\frac{3+\sqrt{57}}{-3+\sqrt{57}}\;\approx\;2.318,
\end{equation}
against the measured $m_{H}^{2}/M_{W}^{2}=2.434$.  The agreement is
$4.8\%$ at the squared ratio.  The same identification at $j=1$ gives
the top mass from $|\Xm(2)|$; the textbook reading has no such
prediction.

\subsection{Identification with the Schur block}

The identification is that the physical $W$ and $Z$ are the positive
Schur eigenvalues of the gauge--scalar block [Eq.~(7) of the Letter].
Writing the quadratic Lagrangian for $(W_{\mu}^{a},\delta H)$ in the
notation of the main Letter,
\begin{equation}
  \mathcal{L}^{(2)}\;=\;\alpha\,\mathrm{Re}\langle a,A_{j}\phi\rangle
  -\beta\,\langle\phi,A_{j}^{\dagger}A_{j}\phi\rangle,
\label{eq:Lalphabeta}
\end{equation}
the cross-coupling coefficient is fixed by \eqref{eq:DmuH} to be
$\alpha=g\sqrt{J}\,v$, and the diagonal scalar mass is the residual
output of the Hosotani potential on $\CP^{1}$.

The single-coupling hypothesis (Proposition~\ref{prop:c1-app}) is the
statement that the Hosotani diagonal is $\beta=\alpha^{2}/g^{2}$, so
that $c=\beta/\alpha^{2}=1$.  Equivalently, the Higgs self-mass is
the rank-one self-energy generated by the same operator that supplies
the mixing.  The structural reason this is reasonable is that
gauge--Higgs unification gives the Higgs only a Hosotani-type mass,
generated by the same gauge interaction \eqref{eq:DmuH}; the open
computation is the explicit one-loop coefficient on the stabilized
$\CP^{1}$ with the Type~IIA flux background, which sets the sign and
the magnitude of $\beta$.

% CONTROL: This is residue (b) of claudeFinal -- the Hosotani-potential
% coefficient.  The PRL companion does not claim to have closed it; it
% claims only that the single-coupling Schur structure is the assumption
% that fixes c=1.  The structural argument is in app:portal below.

\subsection{Contrast with the curl on the cotangent bundle}

A natural alternative would be to identify the off-diagonal operator
with the geometric Hodge curl $\star\dd$ on co-exact one-forms of the
internal $S^{3}=\SU(2)$.  By the Weitzenb\"ock--Lichnerowicz identity
on the unit $S^{3}$, with Ricci tensor $\mathrm{Ric}=2g$,
\begin{equation}
  \Delta_{1}\bigm|_{\rm co\text{-}exact}
  \;=\;\nabla^{*}\nabla+\mathrm{Ric}
  \;=\;\Delta_{0}+1
  \;=\;4J+1\;=\;(2j+1)^{2},
\label{eq:weitz-app}
\end{equation}
so the eigenvalues of the curl on transverse one-forms are
$\pm(2j+1)$.  This is the curvature-shifted normalization.  The
algebraic second-Casimir eigenvalue of Eq.~(8) of the Letter and the
curl eigenvalue \eqref{eq:weitz-app} differ by one unit of curvature:
a consequence of building the off-diagonal from a curvature-bearing
metric operator.  The bare gauge current of gauge--Higgs unification
\eqref{eq:DmuH} delivers the Casimir eigenvalue; an independent
curvature-shifted operator would deliver $4J+1$.  The numerical
distinction between the two normalizations is given in
Sec.~\ref{app:weitz} below.

% ============================================================
\section{Single-coupling Schur completion}
\label{app:portal}
% ============================================================

We prove Proposition~1 of the Letter.  Consider the quadratic action
\eqref{eq:Lalphabeta} for $(a,\phi)$ on the spin-$j$ sector, with
$\alpha,\beta$ real constants, and assume the order-parameter scalar
$\phi$ satisfies the two structural hypotheses

\begin{enumerate}
\item[(P1)] (Flatness) The fluxless mass of $\phi$ vanishes: the bare
quadratic term $m_{0}^{2}|\phi|^{2}$ is absent, so all of $\phi$'s
diagonal mass comes from the gauge sector.
\item[(P2)] (Single Stückelberg coupling) The only coupling between $\phi$ and the
gauge field $a$ is the term $\alpha\,a^{\dagger}A_{j}\phi+\mathrm{h.c.}$
with the same $A_{j}$ that defines the mixing.
\end{enumerate}

\begin{proposition}[Single-coupling Schur completion]
\label{prop:c1-app}
Under (P1) and (P2), integration of the gauge fluctuation in its own
equation of motion generates the diagonal self-energy
\begin{equation}
  M_{\rm s}^{2}\;=\;-\,\alpha^{2}\,A_{j}A_{j}^{\dagger}\,/g^{2},
\label{eq:diag}
\end{equation}
so the reduced fluctuation block on $V_{j}\oplus\mathrm{Im}\,A_{j}$ is
\begin{equation}
  \mu^{2}
  \begin{pmatrix} 0 & A_{j}^{\dagger}\\ A_{j} & -A_{j}A_{j}^{\dagger}\end{pmatrix},
  \qquad
  \mu^{2}=\frac{\alpha^{2}}{g^{2}}\,,
\label{eq:Schur}
\end{equation}
with $c=\beta/\alpha^{2}=1$.
\end{proposition}

\begin{proof}
By (P1), the bare quadratic action for $\phi$ has no diagonal mass.
By (P2), the only quadratic term linking $\phi$ to $a$ is the coupling
$\alpha\,a^{\dagger}A_{j}\phi+\mathrm{h.c.}$.  The gauge equation of
motion is therefore
\begin{equation}
  g^{2}a=\alpha\,A_{j}^{\dagger}\phi.
\label{eq:aeom}
\end{equation}
Substituting \eqref{eq:aeom} back into the action gives the
$\phi$-only quadratic term
\begin{equation}
  -\frac{\alpha^{2}}{g^{2}}\,\langle\phi,A_{j}A_{j}^{\dagger}\phi\rangle,
\label{eq:tracedout}
\end{equation}
which is the rank-one self-energy \eqref{eq:diag}.  Hence
$\beta=\alpha^{2}/g^{2}$ and $c=1$.  Projecting to the image two-channel
subspace and using the Casimir factorization of
Lemma~\ref{lem:cas} reduces \eqref{eq:Schur} to $\mu^{2}\cA(J)$.
\end{proof}

The structural content of (P1)--(P2) is exactly the defining feature
of gauge--Higgs unification: (P1) is the Hosotani flatness of the
Wilson-line modulus, and (P2) is the single $F_{\mu a}$ cross-coupling of
\eqref{eq:DmuH}.  What remains open is the value of the coefficient
$\beta$ from the explicit one-loop Hosotani potential on the
stabilized $\CP^{1}$.  The proposition says that any single-coupling
flat modulus has $\beta=\alpha^{2}/g^{2}$, regardless of the magnitude
of $\alpha$; the magnitude of the off-diagonal squared in $\cA(J)$ is
the bare $\SU(2)$ Casimir $J$, set by Eq.~(8) of the Letter.

% CONTROL: residue (a) -- single-coupling/Schur identification -- is
% claimed here as a proposition.  Its physical content is the
% Hosotani-coefficient sign and magnitude (residue (b)); we do not
% claim to have computed beta from first principles in this PRL.

% ============================================================
\section{The Weitzenb\"ock dichotomy}
\label{app:weitz}
% ============================================================

The bare-Casimir normalization of Eq.~(8) of the Letter and the
curvature-shifted normalization \eqref{eq:weitz-app} produce
distinguishable weak angles.  For a deformation
$\cA_{c}(J)=[[0,\sqrt{J}],[\sqrt{J},-cJ]]$, the positive root is
\begin{equation}
  \Xp(J;c)\;=\;\frac{-cJ+\sqrt{c^{2}J^{2}+4J}}{2},
\label{eq:Xpc}
\end{equation}
and the weak angle is $\sinsq(c)=1-\Xp(\tfrac34;c)/\Xp(2;c)$.
Table~\ref{tab:cscan} records four reference points.

\begin{table}[h]
\centering
\begin{tabular}{lccc}
\toprule
Source & $c_{W}$ at $J=3/4$ & $c_{Z}$ at $J=2$ & $\sinsq$\\
\midrule
Single-coupling Schur (algebraic) & $1$ & $1$ & $0.22310$\\
Weitzenb\"ock-shifted curl & $4/3$ & $9/8$ & $0.267$\\
Conformal weight ($L_{0}/(k+2)$) & $1/2$ & $1/2$ & $0.30$\\
Standard gauge--Higgs (Manton/Fairlie) & --- & --- & $1/4$\\
\bottomrule
\end{tabular}
\caption{The weak angle predicted by four distinct normalizations of
the gauge--scalar mass operator.  Only the bare $\SU(2)$ Casimir of
gauge--Higgs unification reproduces the measured value.  The local
sensitivity at the single-coupling point $c=1$ is $\dd\sinsq/\dd c=-0.140$.}
\label{tab:cscan}
\end{table}

The local rigidity is what makes the relation falsifiable.  A
one-percent error in $c$ produces a shift in $\sinsq$ of $1.4\times
10^{-3}$, an order of magnitude larger than the experimental
uncertainty.  Future improvement of the on-shell measurement to the
$10^{-5}$ level would test the bare Casimir against alternative
near-unity values of $c$.

% ============================================================
\section{Realization in Type IIA on \texorpdfstring{$M_{4}\times\CP^{2}\times\CP^{1}$}{M4 x CP2 x CP1}}
\label{app:realization}
% ============================================================

The operator $\cA(J)$ enters the linearized fluctuation spectrum of a
ten-dimensional Type~IIA flux compactification on $M_{4}\times K_{6}$
with the minimal internal manifold
\begin{equation}
  K_{6}\;=\;\CP^{2}\times\CP^{1}.
\label{eq:K6}
\end{equation}
This is the product of the smallest $\SU(3)$ coset and the smallest
$\SU(2)$ coset, saturating the internal dimension six.  The
$\CP^{1}$ factor supplies the Borel--Weil sectors \eqref{eq:bw-app}.
The $\CP^{2}$ factor carries color and hypercharge.

\subsection{The eleven-dimensional parent as a controlled limit}

The Freund--Rubin family $M^{pqr}$ of homogeneous spaces
\cite{BailinLove1984} admits the circle fibration
$S^{1}_{Z}\to M^{pqr}\to\CP^{2}\times\CP^{1}$, so $K_{6}$ is the
ten-dimensional reduction of $M^{pqr}$.  The eleven-dimensional limit
is well controlled in this sense, and it supplies the linearization
that allows the gauge--scalar block to be read.  The integer labels
$(p,q,r,n)$ play no role in the present construction: the
Bailin--Love coupling ratio is one-to-one with a four-integer family
that is dense around any target, so the precision of any integer
match expresses the density of the family.  This
overdetermination problem is the reason the present route relies on
the algebraic operator $\cA(J)$.  The explicit scan demonstrating the
density of integer matches is recorded in Sec.~\ref{app:scan}.

\subsection{The nine-dimensional electroweak endpoint}

After the single-coupling scalar of Proposition~\ref{prop:c1-app}
develops its vacuum, the weak $\SU(2)$ is broken to the
electromagnetic stabilizer.  Let $P_{T_{3}}$ and $P_{Y}$ denote the
unit frame bundles of the weak weight line and the hypercharge line.
The neutral character is
\begin{equation}
  \chi_{N}(u_{T},u_{Y})\;=\;u_{T}^{-1}u_{Y},
\end{equation}
and the diagonal frame circle is
\begin{equation}
  S^{1}_{Q}\;=\;\ker\chi_{N}\;\subset\;P_{T_{3}}\times P_{Y},
\end{equation}
the line bundle whose first Chern class is
$c_{1}(L_{Y})-c_{1}(L_{T_{3}})$.  The infrared geometry is
$M_{4}\times\CP^{2}\times S^{1}_{Q}$, and the surviving photon is the
zero mode of the diagonal $\U(1)_{Q}$ generator $Q=T_{3}+Y$.  The
photon is exactly massless by \eqref{eq:Qeta}.

The global form of the surviving gauge group is
$(\SU(3)_{c}\times\U(1)_{Q})/\Z_{3}$.  Writing the integer hypercharge
$y=6Y$, the diagonal $\Z_{6}$ generator of the Standard-Model gauge
group is
$\zeta_{6}=(\omega_{3},-1,e^{i\pi/3})$ with $\omega_{3}=e^{2\pi i/3}$;
on a surviving state with color triality $t$, weak weight $T_{3}=m$,
and charge $Q=m+Y$, it acts as
$\omega_{3}^{t}e^{2\pi i Q}$, which reduces to
$\zeta_{3}=(\omega_{3},e^{2\pi i/3})$ on the residual group.  The
descent $\Z_{6}\to\Z_{3}$ is therefore the kinematic content of the
single-coupling symmetry breaking.

\subsection{Chirality and the role of $D=10$}

Chiral fermions are obstructed for any smooth eleven-dimensional
Kaluza--Klein reduction \cite{Witten1981}; the obstruction is removed
in ten dimensions by D-branes wrapping the weak $\CP^{1}$ and
carrying the Standard-Model chirality lattice \cite{Witten1995}.  The
nine-dimensional infrared residue $\SU(3)_{c}\times\U(1)_{Q}$ is
vector-like and admits no chiral obstruction.  The conclusion is that
the physical rung is $D=10$: the operator $\cA(J)$, the chiral
fermions, and the single-coupling Higgs all live on the same ten-dimensional
interface, with the eleven-dimensional parent supplying linearization
and the nine-dimensional residue supplying the vacuum.

% CONTROL: D=10 is rule-1 of AGENT.md.  The geometric chain in this
% subsection is mandatory.

\subsection{The finite Higgs sheet}

The single-coupling scalar of Proposition~\ref{prop:c1-app} admits an
almost-commutative interpretation in the sense of Connes
\cite{ConnesLott1990,ChamseddineConnes2012}.  Writing the
finite-direction Dirac operator
$D_{F}=\bigl(\begin{smallmatrix}0&\Phi\\\Phi^{\dagger}&0\end{smallmatrix}\bigr)$,
the spectral distance to the vacuum is
$d_{\rm sheet}=1/\|\Phi\|$ \cite{MartinettiWulkenhaar2002}.  Identifying
the finite sheet length with the vacuum expectation value through
$\ell_{H}=\sqrt 2/v_{\rm EW}$ ties the electroweak scale to a
finite-geometric datum: the Higgs vacuum amplitude is a length in the
finite direction of the spectral triple.  The negative-branch
prediction $v_{\rm EW}=\sqrt{2}\mu\sqrt{1+\sqrt 3}$ [Eq.~(15) of the
Letter] then reads as a geometric identification: the finite sheet
length and the vacuum amplitude are the same datum.

% ============================================================
\section{The negative branch and the electroweak scale}
\label{app:negative}
% ============================================================

The negative root $\Xm(J)=-J-\Xp(J)$ is an unavoidable consequence of
the indefinite Schur block \eqref{eq:Schur}.  For $J_{1/2}=\tfrac34$
and $J_{1}=2$,
\begin{align}
  \Xm(3/4) &= \tfrac18(-3-\sqrt{57}) = -1.318\dots,\\
  \Xm(2) &= -1-\sqrt 3 = -2.732\dots,
\end{align}
both negative.  Anchoring the overall scale $\mu$ to the $Z$ pole
through $M_{Z}^{2}=\mu^{2}\Xp(2)$ gives
\begin{equation}
  \mu\;=\;\frac{M_{Z}}{\sqrt{\Xp(2)}}
  \;=\;\frac{91.19\,\text{GeV}}{\sqrt{\sqrt 3-1}}
  \;=\;106.6\,\text{GeV},
\label{eq:mu}
\end{equation}
and the $W$ pole follows from $M_{W}=\mu\sqrt{\Xp(3/4)}=80.37\,\text{GeV}$,
in agreement with $M_{W}=80.377(12)\,\text{GeV}$.

The negative-branch eigenvalues set companion scales,
\begin{align}
  \mu\sqrt{|\Xm(3/4)|} &= 122.4\,\text{GeV},
  \label{eq:Higgs}\\
  \mu\sqrt{|\Xm(2)|} &= 176.2\,\text{GeV}.
  \label{eq:top}
\end{align}
The values \eqref{eq:Higgs}--\eqref{eq:top} are within $3\,\text{GeV}$
of the measured Higgs and top masses ($125.25\,\text{GeV}$,
$172.7\,\text{GeV}$), without further parameter fitting.  Identifying
the negative branch of the $j=1$ sector with the electroweak vacuum
amplitude through $v_{\rm EW}=\sqrt 2\mu\sqrt{|\Xm(2)|}$ gives
\begin{equation}
  v_{\rm EW}\;=\;\sqrt 2\,\mu\sqrt{1+\sqrt 3}\;=\;249\,\text{GeV},
\label{eq:vEW-app}
\end{equation}
to be compared with the measured $246.2\,\text{GeV}$.

The identification of $\Xm(J)$ with the order-parameter spectrum is
the geometric content of the radion/Higgs reading.  It is consistent
within the model and agrees with data at the few-GeV level; a
first-principles derivation of the negative-branch identification
requires the explicit Hosotani potential beyond the quadratic order
displayed in Proposition~\ref{prop:c1-app}.

\subsection{Orbital-velocity form and two-description correspondence}

The eigenvalue equation $X^{2}+JX-J=0$ admits a kinematic
interpretation.  Setting $X_{+}=\beta^{2}$ in the orbital-velocity
reading,
\begin{equation}
  \frac{\beta^{2}}{\sqrt{1-\beta^{2}}}\;=\;\sqrt J,
\label{eq:orbit}
\end{equation}
which on squaring reproduces $X^{2}+JX-J=0$.  The gauge-boson masses
are therefore $M_{V}=\mu\beta_{V}$ with $\beta_{W}=0.7541$,
$\beta_{Z}=0.8556$.  The bound $\Xp(J)<1$ that keeps the electroweak
masses below the compactification scale is the relativistic bound
$\beta<1$.

The same eigenvalue equation $X^{2}+JX-J=0$ thus arises from two
unrelated structures: the field-theoretic spin-$j$ fluctuation block
of the Kaluza--Klein reduction, and the kinematic quantization
\eqref{eq:orbit} of a single relativistic mode with internal orbital
momentum $\sqrt{J}$.  Whether the agreement reflects a duality
between the two formulations, or coincidence of the underlying
algebra, is open; a precise field-theory derivation of \eqref{eq:orbit}
from the same Kaluza--Klein action would settle it.

% ============================================================
\section{The 96.5\,GeV prediction and the two-spin tower}
\label{app:tower}
% ============================================================

The same operator at higher half-integer rungs predicts a sequence of
neutral spin-zero states.  The next is
\begin{equation}
  M_{3/2}\;=\;\mu\sqrt{\Xp(15/4)}\;=\;96.5\,\text{GeV}.
\label{eq:96-app}
\end{equation}
A persistent excess in light-Higgs searches at this mass has been
reported by CMS in di-photon final states \cite{CMS96} and by
combinations of LEP and LHC data \cite{Heinemeyer2019}.  The prediction
\eqref{eq:96-app} is conditional on the physical-tower truncation of
$\cA(J)$, in which all $j$ are dynamical.  The alternative finite
truncation at $j\le 1$, natural in an internal geometry of WZW level
two~\cite{Bouwknegt1995}, would forbid the state.

The two-spin tower is the generalization to any pair of internal spins,
\begin{equation}
  \sin^{2}\theta(j_{a},j_{b})\;=\;1-\frac{\Xp(J_{a})}{\Xp(J_{b})}.
\label{eq:tower-app}
\end{equation}
Sample values are tabulated below; the electroweak case is the first
row.

\begin{table}[h]
\centering
\begin{tabular}{lccl}
\toprule
$(j_{a},j_{b})$ & $J_{a},J_{b}$ & $\sin^{2}\theta(j_{a},j_{b})$ & Identification\\
\midrule
$(\half,1)$ & $\tfrac34,\,2$ & $0.22310$ & electroweak\\
$(\half,\tfrac32)$ & $\tfrac34,\,\tfrac{15}4$ & $0.30684$ & prediction\\
$(1,\tfrac32)$ & $2,\,\tfrac{15}4$ & $0.10778$ & prediction\\
$(\half,2)$ & $\tfrac34,\,6$ & $0.34852$ & prediction\\
\bottomrule
\end{tabular}
\caption{The two-spin tower $\sin^{2}\theta(j_{a},j_{b})$ for the four
lowest pairings.  The electroweak entry is forced; the others are
predictions of the same operator for any extended sector with higher
internal spins.}
\label{tab:tower-app}
\end{table}

% ============================================================
\section{The integer-fit overdetermination, recorded for the avoidance of doubt}
\label{app:scan}
% ============================================================

The Bailin--Love coupling ratio for the homogeneous space $M^{pqr}$
gives, for $0<\beta<\half$,
\begin{equation}
  4\beta^{3}-6\beta^{2}+\Bigl(\tfrac94+\frac{q^{2}}{p^{2}}\Bigr)\beta
  -\half\frac{q^{2}}{p^{2}}\;=\;0,
\label{eq:BL}
\end{equation}
and
\begin{equation}
  \frac{g_{Y}^{2}}{g_{2}^{2}}\;=\;
  \frac32\frac{p^{2}n^{2}}{r^{2}}\frac{(3-4\beta)^{2}(1+\beta)}{1-2\beta}.
\label{eq:gYg2}
\end{equation}
A modest integer scan over $p,q\le 30$, $r\le 300$, $n\le 3$ gives 138
four-tuples within $10^{-3}$ of the target ratio $0.287169$, 12 within
$10^{-4}$, and 2 within $10^{-5}$.  The achievable values are dense
near the target; the six-digit agreement of the tuple $(12,14,91,1)$
expresses the density of the family.

This overdetermination is the reason the present construction relies
on $\cA(J)$ as the weak-angle source.  No integer label is used in any
derivation.  The labels appear only when one wishes to embed the
$\CP^{2}\times\CP^{1}$ geometry inside a Freund--Rubin parent
$M^{pqr}$.  There they are constrained by global-form data: the
$\Z_{6}$ quotient of the Standard-Model gauge group, the spin$^{c}$
structure on $\CP^{2}$, and the period of the hypercharge line bundle.
The constraints leave the labels underdetermined at selection
precision.

% CONTROL: this appendix is the explicit warning that the integer-fit
% route is not used.  Including it here protects the PRL against the
% objection that any K-K construction with M^{pqr} input is implicitly
% fitting the angle through (p,q,r,n).

% ============================================================
\section{Pole-mass vs. running-coupling reading of \texorpdfstring{$\sinsq$}{sin2thetaW}}
\label{app:scheme}
% ============================================================

The predicted $\sinsq=0.22310$ is the on-shell ratio
$1-M_{W}^{2}/M_{Z}^{2}$.  The relevant data point is the on-shell
value $0.22305(20)$; the $\overline{\rm MS}$ coupling angle
$\hat s_{Z}^{2}\simeq 0.23122$ at $\mu=M_{Z}$ is a different
observable.  A
geometric construction of \emph{couplings} would more naturally yield
an $\overline{\rm MS}$-like quantity; that $\cA(J)$ lands on the
on-shell mass ratio reflects that the matrix is a \emph{mass}
operator, generated by the rank-one Schur completion of a single
Stückelberg coupling.

The complex pole of the gauge-boson propagator is
$s_{V}=M_{V}^{2}-iM_{V}\Gamma_{V}$, and two on-shell conventions
circulate: the running mass $\overline{M}_{V}$ and the Breit--Wigner
on-shell mass $M_{V}$, differing by
$M_{V}-\overline{M}_{V}\simeq\Gamma_{V}^{2}/(2M_{V})$.  This is
$34\,\text{MeV}$ for the $Z$ and $27\,\text{MeV}$ for the $W$,
translating into a few-times-$10^{-4}$ ambiguity in the ratio, which
is the size of the residual between $0.22310$ and $0.22305$.  The
agreement is therefore within the scheme ambiguity.

% ============================================================
\section{Numerical summary}
\label{app:numbers}
% ============================================================

\begin{table}[h]
\centering
\begin{tabular}{lcc}
\toprule
Quantity & Value (this model) & Comparison\\
\midrule
$\Xp(3/4)$ & $0.568729$ & ---\\
$\Xp(2)$ & $0.732051=\sqrt 3-1$ & ---\\
$\sinsq=1-\Xp(3/4)/\Xp(2)$ & $0.223101$ & $0.22305(20)$\,~\cite{PDG2024}\\
$M_{W}^{2}/M_{Z}^{2}$ & $0.776898$ & $0.77687(2)$\\
$\rho_{\rm cust}$ & $1$ & $1.00031(19)$\\
$\mu=M_{Z}/\sqrt{\Xp(2)}$ & $106.6\,\text{GeV}$ & ---\\
$M_{W}=\mu\sqrt{\Xp(3/4)}$ & $80.37\,\text{GeV}$ & $80.377(12)$\\
$\mu\sqrt{|\Xm(3/4)|}$ & $122.4\,\text{GeV}$ & $m_{H}=125.25$\\
$\mu\sqrt{|\Xm(2)|}$ & $176.2\,\text{GeV}$ & $m_{t}=172.7$\\
$M_{3/2}=\mu\sqrt{\Xp(15/4)}$ & $96.5\,\text{GeV}$ & prediction\\
$v_{\rm EW}=\sqrt 2\mu\sqrt{1+\sqrt 3}$ & $249\,\text{GeV}$ & $246.2$\\
$\dd\sinsq/\dd c$ at $c=1$ & $-0.140$ & rigidity slope\\
\bottomrule
\end{tabular}
\caption{Numerical content of the construction.  Experimental values
from \cite{PDG2024}.  The first three rows are the load-bearing
predictions; rows $4$--$5$ are immediate consequences; rows $6$--$10$
follow from the negative-branch identification.}
\label{tab:numerics}
\end{table}

% ============================================================
\section{Open computations}
\label{app:open}
% ============================================================

Three computations sharpen the construction beyond its present state.
First, the Hosotani potential on the stabilized $\CP^{1}$ with the
Type~IIA flux background determines the coefficient $\beta$ in
\eqref{eq:Lalphabeta} and so the magnitude of the negative diagonal;
the single-coupling Schur completion of Proposition~\ref{prop:c1-app}
fixes the ratio $c=1$ but not the absolute magnitude.  Second, the
explicit linearized IIA reduction on $M_{4}\times\CP^{2}\times\CP^{1}$
with stabilized moduli verifies that the order-parameter scalar is
indeed the single-coupling flat modulus assumed by (P1)--(P2).  Third,
the spin$^{c}$-extended fermion content on $\CP^{2}$ supplies the
chiral matter of the Standard Model, including the hypercharge
assignment $Y=-\half$ on the right-handed singlet that closes the
determinant-one lift of the gauge group.

The Hosotani-potential coefficient, the linearized IIA reduction on
the stabilized internal manifold, and the spin$^{c}$-extended fermion
content are computations whose completion sharpens the construction
beyond its present state.  The present Letter establishes the
operator-level content of the electroweak mixing angle and the
on-shell mass ratio, the structural reason for the bare-Casimir
normalization through gauge--Higgs unification, the single-coupling
Schur completion that fixes $c=1$, the Weitzenb\"ock rigidity of the
prediction, and the $96.5\,\text{GeV}$ scalar that follows from the
same operator.

% ============================================================
% Bibliography.  Each entry carries a `% CITED IN' comment indicating
% the supplemental section where it is invoked.  Entries cross-listed
% with main.tex carry the same key.
% ============================================================
\begin{thebibliography}{99}

% CITED IN: numerical-summary table (Sec.~VIII); on-shell scheme paragraph
% in Sec.~IX (comparison to 0.22305(20)).
\bibitem{PDG2024}
S.~Navas \emph{et~al.} (Particle Data Group), Phys.\ Rev.\ D
\textbf{110}, 030001 (2024).

% CITED IN: realization section (Sec.~V), chirality paragraph (Witten
% no-go for smooth 11D Kaluza-Klein).
\bibitem{Witten1981}
E.~Witten, Nucl.\ Phys.\ B \textbf{186}, 412 (1981).

% CITED IN: realization section (Sec.~V), chirality paragraph (D-branes
% in ten dimensions supply chiral matter via the M-theory web).
\bibitem{Witten1995}
E.~Witten, ``String theory dynamics in various dimensions,'' Nucl.\ Phys.\
B \textbf{443}, 85 (1995); arXiv:hep-th/9503124.

% CITED IN: realization section (Sec.~V), M^{pqr} circle reduction giving
% K_6 = CP^2 x CP^1.
\bibitem{BailinLove1984}
D.~Bailin and A.~Love, Phys.\ Lett.\ B \textbf{144}, 359 (1984).

% CITED IN: comparison-with-textbook subsection (Sec.~II.D).  Original
% gauge-Higgs unification reference giving sin^2 theta_W = 1/4 on SU(3)/U(2).
\bibitem{Manton1979}
N.~S.~Manton, Nucl.\ Phys.\ B \textbf{158}, 141 (1979).

% CITED IN: comparison-with-textbook subsection (Sec.~II.D).  Original
% Higgs-from-internal-gauge-field paper for SU(3) on S^2.
\bibitem{Fairlie1979}
D.~B.~Fairlie, Phys.\ Lett.\ B \textbf{82}, 97 (1979).

% CITED IN: finite-Higgs-sheet paragraph in realization section (Sec.~V),
% almost-commutative interpretation of the order parameter.
\bibitem{ConnesLott1990}
A.~Connes and J.~Lott, Nucl.\ Phys.\ B Proc.\ Suppl.\ \textbf{18B}, 29
(1990).

% CITED IN: finite-Higgs-sheet paragraph in realization section (Sec.~V),
% spectral-action context for the same.
\bibitem{ChamseddineConnes2012}
A.~H.~Chamseddine and A.~Connes, ``Resilience of the spectral standard
model,'' JHEP \textbf{09}, 104 (2012); arXiv:1208.1030.

% CITED IN: finite-Higgs-sheet paragraph in realization section (Sec.~V),
% scalar-fluctuation distance formula d_sheet = 1/||Phi||.
\bibitem{MartinettiWulkenhaar2002}
P.~Martinetti and R.~Wulkenhaar, J.~Math.~Phys.~\textbf{43}, 182
(2002); arXiv:hep-th/0104108.

% CITED IN: 96.5 GeV / two-spin tower section (Sec.~VII), level-two
% truncation of integrable SU(2)_k sectors.
% [verified arXiv:hep-th/9412108; "Spinon basis for higher level SU(2)
% WZW models," Bouwknegt, Ludwig, Schoutens.]
\bibitem{Bouwknegt1995}
P.~Bouwknegt, A.~W.~W.~Ludwig, and K.~Schoutens, ``Spinon basis for
higher level SU(2) WZW models,'' arXiv:hep-th/9412108.

% CITED IN: 96.5 GeV / two-spin tower section (Sec.~VII), CMS diphoton
% excess at 95.4 GeV.
% [verified arXiv:2405.18149; corrected from earlier draft that cited the
% tau-tau analysis arXiv:2208.02717.]
\bibitem{CMS96}
The CMS Collaboration, ``Search for a standard model-like Higgs boson in
the mass range between 70 and 110 GeV in the diphoton final state in
proton-proton collisions at $\sqrt{s}=13$ TeV,'' arXiv:2405.18149.

% CITED IN: 96.5 GeV / two-spin tower section (Sec.~VII), LEP+LHC
% combination interpretation as a 96 GeV Higgs.
% [verified arXiv:1812.05864; "A Higgs Boson at 96 GeV?!" Heinemeyer and
% Stefaniak, PoS CHARGED2018, 016 (2019).]
\bibitem{Heinemeyer2019}
S.~Heinemeyer and T.~Stefaniak, ``A Higgs Boson at 96 GeV?!,'' PoS
\textbf{CHARGED2018}, 016 (2019); arXiv:1812.05864.

\end{thebibliography}

\end{document}
